\documentclass{article}
\usepackage{mathtools} 
\usepackage{fontspec}
\usepackage[UTF8]{ctex}
\usepackage{amsthm}
\usepackage{mdframed}
\usepackage{xcolor}
\usepackage{amssymb}
\usepackage{amsmath}


% 定义新的带灰色背景的说明环境 zremark
\newmdtheoremenv[
  backgroundcolor=gray!10,
  % 边框与背景一致，边框线会消失
  linecolor=gray!10
]{zremark}{说明}

% 通用矩阵命令: \flexmatrix{矩阵名}{元素符号}{行数}{列数}
\newcommand{\flexmatrix}[4]{
  \[
  #1 = \begin{pmatrix}
    #2_{11}     & #2_{12}     & \cdots & #2_{1#4}   \\
    #2_{21}     & #2_{22}     & \cdots & #2_{2#4}   \\
    \vdots      & \vdots      & \ddots & \vdots     \\
    #2_{#31}    & #2_{#32}    & \cdots & #2_{#3#4}
  \end{pmatrix}
  \]
}

% 简化版命令（默认矩阵名为A，元素符号为a）: \quickmatrix{行数}{列数}
\newcommand{\quickmatrix}[2]{\flexmatrix{A}{a}{#1}{#2}}

\begin{document}
\title{注释 13.1}
\author{张志聪}
\maketitle

\begin{zremark}
  证明：定理13.2的推论。
\end{zremark}
\textbf{证明：}
若不一致收敛于$f$，即
\begin{align*}
  \lim\limits_{n \to +\infty} \sup\limits_{x \in D} |f_n(x) - f(x)| \neq 0
\end{align*}
由极限的定义，存在$\epsilon_0 > 0$，对任意$N > 0$，都存在$n > N$，使得
\begin{align*}
  \left|\sup\limits_{x \in D} |f_n(x) - f(x)| - 0 \right| & \geq \epsilon_0 \\
  \sup\limits_{x \in D} |f_n(x) - f(x)|                   & \geq \epsilon_0
\end{align*}
由上确界的定义，存在$x_n \in D$，使得
\begin{align*}
  |f_n(x) - f(x)| > \frac{1}{2} \epsilon_0
\end{align*}
于是得到数列$\{x_n\}$，使得
\begin{align*}
  |f_n(x_n) - f(x_n)| \geq \frac{1}{2} \epsilon_0
\end{align*}
如果极限存在，则
\begin{align*}
  \lim\limits_{n \to +\infty} |f_n(x_n) - f(x_n)| \geq \frac{1}{2} \epsilon_0
\end{align*}
另外，极限不存在，也满足推论中的描述。

\begin{zremark}
  步骤说明：
  \begin{align*}
    \sup\limits_{x \in (-1, 1)} \left| \frac{x^n}{x - 1} \right|
    \geq \left| \frac{\left( \frac{n}{n+1} \right)^n}{1 - \frac{n}{n + 1}} \right|
  \end{align*}
  其中$n$是正整数。
\end{zremark}
\textbf{证明：}
因为
\begin{align*}
  -1 < x = \frac{n}{n+1} < 1
\end{align*}
即$x \in (-1, 1)$，由上确界的定义可得以上不等式成立。

\end{document}